% ── SECTION: COHERENCE-BUDGET ALLOCATION (~v1.0 new content) ────────────────
%   v1.0 contribution: the allocation theorem, scheduling algorithm,
%   and multi-layer fidelity benchmark.

\section{Coherence-Budget Allocation for Tunable Exchange Schedules}
\label{sec:alloc}

The per-gate fidelity comparison of Sec.~\ref{sec:sim} establishes
that, on hardware with directly calibratable Heisenberg exchange,
$\URA(\theta)$ at $\theta < \pi/4$ accumulates less coherence-limited
error than a full maximally entangling primitive at equal entangling
capability requirement.
That observation, however, says nothing about how to \emph{allocate}
entanglement across the gate layers of a circuit whose total
entanglement target is fixed.
This section turns the per-gate fidelity statement into a layerwise
optimization: given a coherence budget and an entanglement target,
how should a noise-aware compiler distribute exchange angles across
the circuit?

\subsection{Relational Coherence Cost}
\label{subsec:cost}

I define the per-layer cost of a tunable exchange gate
$\URA(\theta_\ell)$ as the local pointer-basis coherence consumed by
the gate:
%
\begin{equation}
\label{eq:layer-cost}
C_{\RCF}^{(\ell)} \;\equiv\; \bar\alpha_\ell^{\rm pre} - \bar\alpha_\ell^{\rm post}
\;=\; \bar\alpha_\ell^{\rm pre}\bigl(1 - e^{-\gamma_\ell \kappa\, \theta_\ell}\bigr),
\end{equation}
%
where $\gamma_\ell$ is the local pointer-basis dephasing rate active
during layer $\ell$ (which may vary across layers due to hardware
heterogeneity), $t_\ell = \kappa\, \theta_\ell$ is the exchange-pulse
duration on the substrate~\cite{Krantz2019}, and the exponential
form follows from
Proposition~\ref{prop:determination} in its uniform-rate Markovian
regime.
The total circuit cost is
%
\begin{equation}
\label{eq:circuit-cost}
C_{\RCF}(\mathcal{U}) \;=\; \sum_{\ell=1}^{L} C_{\RCF}^{(\ell)}.
\end{equation}
%
To leading order in $\gamma_\ell \kappa\, \theta_\ell \ll 1$
(the small-cost regime in which Eq.~\eqref{eq:peff}'s
gate-time-proportional model itself is justified, Sec.~\ref{sec:sim}),
%
\begin{equation}
\label{eq:linear-cost}
C_{\RCF}(\mathcal{U}) \;\approx\;
  \sum_{\ell=1}^{L} \tilde\gamma_\ell\, \theta_\ell,
  \qquad \tilde\gamma_\ell \equiv \bar\alpha_\ell^{\rm pre}\, \gamma_\ell\, \kappa.
\end{equation}
%
The effective rate $\tilde\gamma_\ell$ folds the substrate's exchange
time-constant and the layer's residual pre-gate coherence into a
single per-layer cost coefficient.

\subsection{Proposition~V: Exchange-Angle Coherence Cost}
\label{subsec:prop5}

The cost per layer admits a sharp closed form for the
directly-calibrated exchange-gate case treated throughout this paper.

\begin{proposition}[Exchange-angle coherence cost, Markovian regime]
\label{prop:exchange-cost}
Under uniform Markovian pointer-basis dephasing at rate $\gamma$ and the
direct exchange-gate duration of Sec.~\ref{sec:sim},
$t(\theta) = t_{\pi/4}\,\theta/(\pi/4)$, the layerwise scalar-affinity loss induced by
$\URA(\theta)$ on a layer with input affinity $\bar\alpha_{\rm in}$ is
%
\begin{equation}
\label{eq:exchange-cost}
\Delta\bar\alpha(\theta)
  \;=\; \bar\alpha_{\rm in}\!\left[\,
        1 - \exp\!\left(-\gamma\,t_{\pi/4}\,\frac{\theta}{\pi/4}\right)\right].
\end{equation}
%
In the small-error regime $\gamma\, t_{\pi/4}\,\theta/(\pi/4) \ll 1$,
%
\begin{equation}
\label{eq:exchange-cost-linear}
\Delta\bar\alpha(\theta)
  \;\approx\; \gamma\,t_{\pi/4}\,\frac{\theta}{\pi/4}\;\bar\alpha_{\rm in}.
\end{equation}
%
The per-layer coherence cost is therefore first-order linear in the
exchange angle $\theta$.
\end{proposition}

\begin{proof}
By Proposition~\ref{prop:determination} (Markovian regime), the scalar
affinity decays as
$\bar\alpha(t) = \bar\alpha_{\rm in}\,e^{-\gamma t}$.
Substituting $t = t_{\pi/4}\,\theta/(\pi/4)$ gives
$\bar\alpha^{\rm post} = \bar\alpha_{\rm in}\,\exp(-\gamma\,t_{\pi/4}\,\theta/(\pi/4))$,
and $\Delta\bar\alpha \equiv \bar\alpha_{\rm in} - \bar\alpha^{\rm post}$
gives Eq.~\eqref{eq:exchange-cost}.
The linearization Eq.~\eqref{eq:exchange-cost-linear} follows from
$1 - e^{-x} = x + O(x^2)$. $\square$
\end{proof}

\begin{remark}[Linear cost matches Eq.~\eqref{eq:linear-cost}]
The leading-order form Eq.~\eqref{eq:exchange-cost-linear} reproduces
the per-layer term in Eq.~\eqref{eq:linear-cost} with effective rate
$\tilde\gamma_\ell = \bar\alpha_\ell^{\rm pre}\,\gamma_\ell\,t_{\pi/4}/(\pi/4)$;
the constant $\kappa = t_{\pi/4}/(\pi/4)$ is the substrate's
exchange-time scale per unit angle.
The allocation problem of Sec.~\ref{subsec:problem} below is therefore
the layer-aggregated form of
Proposition~\ref{prop:exchange-cost}.
\end{remark}

\subsection{Corollary: Entanglement Efficiency}
\label{subsec:eff}

Combining Proposition~\ref{prop:exchange-cost} with the
entanglement-capability formula
$\mathcal{E}(\theta) = |\sin 2\theta|$ (Sec.~\ref{subsec:entcap})
yields the per-gate entanglement-per-coherence-cost ratio.

\begin{corollary}[Entanglement efficiency]
\label{cor:eta}
In the small-error regime of
Proposition~\ref{prop:exchange-cost}, the entanglement
generated per unit of pointer-basis coherence spent on a single
exchange gate is
%
\begin{equation}
\label{eq:eta}
\eta(\theta) \;\equiv\;
  \frac{\mathcal{E}(\theta)}{\Delta\bar\alpha(\theta)}
  \;\propto\; \frac{\sin 2\theta}{\theta},
\end{equation}
%
suppressing the constant factor
$\bigl(\gamma\,t_{\pi/4}\,\bar\alpha_{\rm in}/(\pi/4)\bigr)^{-1}$.
The ratio $\sin(2\theta)/\theta$ is strictly decreasing on
$(0, \pi/4]$, with $\lim_{\theta\to 0^+}\sin(2\theta)/\theta = 2$
and $\sin(\pi/2)/(\pi/4) = 4/\pi \approx 1.273$.
\end{corollary}

\begin{remark}[Efficiency vs.\ capacity]
Corollary~\ref{cor:eta} formalizes the tension exploited by the
allocation theorem.
Small-$\theta$ gates are most \emph{efficient} per unit coherence
spent (high $\eta$), but produce small absolute entanglement
$\mathcal{E}(\theta) \approx 2\theta$ per gate.
Large-$\theta$ gates are less efficient ($\eta \to 4/\pi$ at the
$\sqrt{\textsc{swap}}$ point) but achieve more capacity per gate.
The Coherence-Budget Allocation Theorem
(Sec.~\ref{subsec:thm-alloc}) resolves this tension by distributing
entanglement across layers when $L > 1$ allows, and concentrates it
when noise heterogeneity forces a tradeoff.
\end{remark}

\subsection{The Allocation Problem}
\label{subsec:problem}

For a circuit with required total entangling capability
$E_{\rm target}$ accumulated across $L$ exchange layers, each
contributing $\mathcal{E}(\theta_\ell) = \sin(2\theta_\ell)$
(Sec.~\ref{subsec:entcap}; I restrict to the entangling half
$\theta_\ell \in [0, \pi/4]$ on which $\sin(2\theta_\ell)$ is
increasing), the coherence-optimal schedule solves
%
\begin{equation}
\label{eq:alloc-problem}
\begin{aligned}
\text{minimize}\quad & C_{\RCF}(\boldsymbol\theta) \;=\;
   \sum_{\ell=1}^L \tilde\gamma_\ell\, \theta_\ell \\
\text{subject to}\quad & \sum_{\ell=1}^L \sin(2\theta_\ell) \;\geq\; E_{\rm target}, \\
& 0 \leq \theta_\ell \leq \pi/4, \qquad \ell = 1,\dots, L.
\end{aligned}
\end{equation}
%
The objective is linear (hence convex) and the constraint set is
convex: $\sum_\ell \sin(2\theta_\ell)$ is concave on
$[0,\pi/4]^L$ (second derivative $-4\sin 2\theta_\ell \leq 0$
component-wise), so the superlevel set
$\{\boldsymbol\theta : \sum_\ell \sin(2\theta_\ell) \geq E_{\rm target}\}$
is convex.
Minimizing a linear functional on a convex set yields a unique
minimizer (interior) or a minimizing face (boundary).

\subsection{The Coherence-Budget Allocation Theorem}
\label{subsec:thm-alloc}

\begin{theorem}[Coherence-Budget Allocation]
\label{thm:alloc}
Let $\tilde\gamma_1,\dots,\tilde\gamma_L > 0$ be the effective
per-layer dephasing rates of Eq.~\eqref{eq:linear-cost} and let
$E_{\rm target} \in (0, L)$.
The unique solution of the convex program
Eq.~\eqref{eq:alloc-problem} is given layerwise by
%
\begin{equation}
\label{eq:alloc-solution}
\theta_\ell^{\!*} \;=\;
  \frac{1}{2}\,\arccos\!\Bigl(
    \min\!\Bigl(1,\;\tfrac{\tilde\gamma_\ell}{2\lambda^{\!*}}\Bigr)
  \Bigr),
  \qquad \ell = 1,\dots,L,
\end{equation}
%
where the shadow price $\lambda^{\!*} > 0$ is the unique
solution of the budget constraint
%
\begin{equation}
\label{eq:lambda-eq}
\sum_{\ell=1}^L \sin\!\bigl(2\theta_\ell^{\!*}(\lambda^{\!*})\bigr)
\;=\; E_{\rm target}.
\end{equation}
\end{theorem}

\begin{proof}
The Lagrangian is
$\mathcal{L}(\boldsymbol\theta,\lambda,\boldsymbol\mu,\boldsymbol\nu)
= \sum_\ell \tilde\gamma_\ell\,\theta_\ell
- \lambda\!\bigl(\sum_\ell \sin 2\theta_\ell - E_{\rm target}\bigr)
- \sum_\ell \mu_\ell\,\theta_\ell
+ \sum_\ell \nu_\ell\,(\theta_\ell - \tfrac{\pi}{4})$
with $\lambda,\mu_\ell,\nu_\ell \geq 0$.
Stationarity $\partial\mathcal{L}/\partial\theta_\ell = 0$ gives
%
\begin{equation*}
\tilde\gamma_\ell - 2\lambda \cos 2\theta_\ell - \mu_\ell + \nu_\ell = 0.
\end{equation*}
%
For an interior layer ($0 < \theta_\ell^* < \pi/4$, complementary
slackness $\mu_\ell = \nu_\ell = 0$) this reduces to
$\cos 2\theta_\ell^* = \tilde\gamma_\ell/(2\lambda)$,
giving Eq.~\eqref{eq:alloc-solution} when the argument is in $[0,1]$.
For the lower boundary $\theta_\ell^* = 0$ ($\nu_\ell = 0$,
$\mu_\ell \geq 0$), stationarity requires
$\tilde\gamma_\ell - 2\lambda \geq 0$, i.e., $\tilde\gamma_\ell \geq 2\lambda$,
which is exactly the condition $\tilde\gamma_\ell/(2\lambda) \geq 1$
in Eq.~\eqref{eq:alloc-solution} (the $\min(\cdot)$ clip).
The upper boundary $\theta_\ell^* = \pi/4$ requires $\tilde\gamma_\ell = 0$,
excluded by hypothesis.

Existence and uniqueness of $\lambda^{\!*}$:
$\sin 2\theta_\ell^*(\lambda)$ is monotone non-decreasing and
continuous in $\lambda$ on $(0,\infty)$, with limits $0$ as $\lambda\to 0^+$
and $1$ as $\lambda\to\infty$.
Hence $\Sigma(\lambda) \equiv \sum_\ell \sin 2\theta_\ell^*(\lambda)$
is continuous and monotone, mapping $(0,\infty)$ onto $(0, L)$;
$\Sigma(\lambda^{\!*}) = E_{\rm target} \in (0,L)$ has a unique
solution.
Global optimality follows from the convexity of the program. $\square$
\end{proof}

\subsection{Three Falsifiable Predictions}
\label{subsec:predictions}

Three predictions follow immediately from Eq.~\eqref{eq:alloc-solution}
and are testable in numerical simulation.

\textbf{P1 (Distribution beats concentration, uniform-$\tilde\gamma$).}
For $\tilde\gamma_\ell = \tilde\gamma$ uniform, the optimal schedule
and its total cost are
\[
\theta_\ell^* = \tfrac{1}{2}\arcsin(E_{\rm target}/L)\ \text{on all layers},
\qquad
C^* = \tilde\gamma\,L\cdot \tfrac{1}{2}\arcsin(E_{\rm target}/L).
\]
Since $\arcsin$ is strictly convex on $[0,1]$ with $\arcsin 0 = 0$,
the ratio $\arcsin(x)/x$ is strictly increasing on $(0,1]$; setting
$x = E_{\rm target}/L < E_{\rm target}$ gives
$L \arcsin(E_{\rm target}/L) < \arcsin(E_{\rm target})$ for $L > 1$
and $E_{\rm target} \in (0,1]$;
hence the optimal distributed schedule strictly beats the cost of
implementing $E_{\rm target}$ in a single partial entangler.

\textbf{P2 (Quieter qubits receive more entanglement).}
For heterogeneous $\tilde\gamma_\ell$, layers with smaller
$\tilde\gamma_\ell$ have smaller $\tilde\gamma_\ell/(2\lambda^{\!*})$,
hence smaller $\cos 2\theta_\ell^*$, hence larger $\theta_\ell^*$.
The optimal compiler routes entangling angle toward quieter substrates.

\textbf{P3 (Skip threshold for noisy qubits).}
Layers with $\tilde\gamma_\ell \geq 2\lambda^{\!*}$ are clipped to
$\theta_\ell^* = 0$: the optimum spends no coherence on layers whose
local cost exceeds the shadow price of the global budget constraint,
provided quieter alternatives suffice to meet $E_{\rm target}$.

P1 is a no-heterogeneity baseline statement;
P2 and P3 are the practically distinguishing predictions on
realistic hardware whose dephasing rates vary across qubit pairs.

\subsection{Scheduling Algorithm}
\label{subsec:algorithm}

The shadow-price equation Eq.~\eqref{eq:lambda-eq} is a monotone
one-dimensional root-finding problem and admits a Brent-method
solution in milliseconds for $L \lesssim 10^3$.
The full algorithm is:

\begin{enumerate}
  \item Input: layer rates $\{\tilde\gamma_\ell\}_{\ell=1}^L$,
        entanglement target $E_{\rm target}$.
  \item Bracket $\lambda$: find $[\lambda_{\rm lo}, \lambda_{\rm hi}]$
        such that $\Sigma(\lambda_{\rm lo}) < E_{\rm target} < \Sigma(\lambda_{\rm hi})$.
  \item Solve $\Sigma(\lambda^{\!*}) = E_{\rm target}$ by Brent's method.
  \item Return $\theta_\ell^* = \tfrac{1}{2}\arccos(\min(1,\tilde\gamma_\ell/(2\lambda^{\!*})))$.
\end{enumerate}

The reference implementation (\texttt{figures/\allowbreak rcb\_\allowbreak optimizer.py})
runs in $O(L \log(1/\epsilon))$ time for tolerance $\epsilon$ and is
included with this paper.

\subsection{Multi-Layer Benchmark}
\label{subsec:benchmark-alloc}

I test Theorem~\ref{thm:alloc} against two naive scheduling
baselines on a multi-layer 2-qubit circuit with per-layer independent
dephasing.

\textbf{Setup.}
A circuit of $L = 8$ layers applies $\URA(\theta_\ell)$ followed by
independent per-qubit pointer-basis dephasing at rate $\gamma_\ell$
for duration $t_\ell = \kappa\,\theta_\ell$, starting from
$\rho_0 = \ket{+0}\!\bra{+0}$.
For each entanglement target $E_{\rm target} \in [0.25, 3.5]$, I
compare three strategies:
\begin{itemize}
  \item \textbf{RCB-optimal} schedule per Theorem~\ref{thm:alloc};
  \item \textbf{Uniform $\theta$}: $\theta_\ell = \tfrac{1}{2}\arcsin(E_{\rm target}/L)$
        on all layers (the optimal schedule when $\tilde\gamma_\ell$ is uniform);
  \item \textbf{Concentrated}: place full entanglers
        $\theta_\ell = \pi/4$ on the $\lfloor E_{\rm target}\rfloor$
        quietest layers, residual on the next-quietest.
\end{itemize}
For each strategy I compute the final reduced-state fidelity
$F = \mathrm{Tr}[\sqrt{\sqrt{\rho_{\rm noisy}}\,\rho_{\rm ideal}\,\sqrt{\rho_{\rm noisy}}}]^2$
against the noiseless target unitary.

\textbf{Results.}
Figure~\ref{fig:rcb-sweep} shows fidelity (top row) and coherence cost
(bottom row) as $E_{\rm target}$ varies, for three $\gamma$ regimes:
uniform low ($\gamma = 0.05$), uniform high ($\gamma = 0.15$), and
heterogeneous ($\gamma \in [0.02, 0.30]$).
The RCB-optimal schedule meets or exceeds both baselines for every
$E_{\rm target}$ tested, with the largest fidelity advantage in the
high-$\gamma$ regime (e.g., $\Delta F = +0.040$ over concentrated at
$E_{\rm target} = 1$, $\gamma = 0.15$) and a substantial advantage
over uniform-$\theta$ in the heterogeneous regime
($\Delta F = +0.025$ at $E_{\rm target} = 1$).
At representative targets:

\begin{center}
\small
\begin{tabular}{lccc}
\toprule
Regime & $\Delta F_{\rm opt - conc}$ at $E\!=\!1$ & at $E\!=\!1.5$ & at $E\!=\!2$ \\
\midrule
Uniform low ($\gamma\!=\!0.05$) & $+0.014$ & $+0.014$ & $+0.013$ \\
Uniform high ($\gamma\!=\!0.15$) & $\mathbf{+0.040}$ & $+0.037$ & $+0.033$ \\
Heterogeneous & $+0.005$ & $+0.004$ & $+0.001$ \\
\bottomrule
\end{tabular}
\end{center}

\begin{figure*}[htbp]
\centering
\includegraphics[width=0.9\textwidth]{fig3_rcb_sweep}
\caption{
  Coherence-Budget Allocation benchmark on an $L\!=\!8$ layer
  two-qubit circuit with per-layer independent pointer-basis dephasing.
  \textbf{Top row:} final-state fidelity vs noiseless target, as
  a function of entanglement target $E_{\rm target}$, for the RCB-optimal
  schedule (Eq.~\eqref{eq:alloc-solution}, solid),
  uniform-$\theta$ baseline (dashed), and concentrated full-entangler
  baseline (dotted).
  \textbf{Bottom row:} coherence cost
  $\sum_\ell \tilde\gamma_\ell\,\theta_\ell$ for the same schedules.
  \textbf{Columns:} uniform low noise
  ($\gamma\!=\!0.05$), uniform high noise ($\gamma\!=\!0.15$), and
  heterogeneous noise ($\gamma \in [0.02, 0.30]$).
  The RCB-optimal schedule equals or exceeds both baselines in
  fidelity across all $E_{\rm target}$; the largest advantage appears
  in the high-$\gamma$ regime (Prediction~P1) and in the
  heterogeneous regime against uniform $\theta$ (Prediction~P2).
  Generated by \texttt{figures/\allowbreak benchmark\_\allowbreak sweep.py}.
}
\label{fig:rcb-sweep}
\end{figure*}

P1 is visible in the uniform-$\gamma$ columns: distribution always
beats concentration.
P2 is visible in the heterogeneous column: RCB-optimal beats
uniform-$\theta$ by routing entanglement onto the
$\gamma\!=\!0.02$ layers and zero onto the $\gamma\!=\!0.30$ layer
(P3 in action).
The benchmark validates Theorem~\ref{thm:alloc} as a per-circuit,
not merely per-gate, statement of the coherence-budget framing.

\subsection{Layer-Varying Affinity}
\label{subsec:alpha-varying}

The benchmark above runs every layer at the same pre-gate affinity, and
in that regime $\bar\alpha$ drops out of the optimum: a common factor in
the effective rates $\tilde\gamma_\ell =
\bar\alpha_\ell^{\rm pre}\gamma_\ell\kappa$ is absorbed by the shadow
price (Sec.~\ref{subsec:limitations}, fifth limitation).
The \RCF~cost and a cost built from the dephasing rates alone differ only
when $\bar\alpha_\ell^{\rm pre}$ varies across layers, and this
subsection tests that regime (Fig.~\ref{fig:alpha-varying},
\texttt{figures/\allowbreak benchmark\_\allowbreak alpha\_\allowbreak varying.py}).
The circuits are the two-qubit, $L = 8$ circuits of
Sec.~\ref{subsec:benchmark-alloc}, with two additions that make
$\bar\alpha_\ell^{\rm pre}$ differ between layers: fixed single-qubit
rotations before each exchange layer, and idle dephasing between
layers that is common to every schedule.
The \RCF\ schedule takes $\bar\alpha_\ell^{\rm pre}$ from the noisy
trajectory of the rate-only schedule and then applies
Theorem~\ref{thm:alloc} once.
A self-consistent version, in which the profile is recomputed from the
schedule it produces, does not converge: the linear program moves all
of the entangling weight onto the lowest-affinity layers, which changes
the profile, and the iteration oscillates.
As a reference I maximize the simulated fidelity numerically under the
same constraint $\sum_\ell \sin 2\theta_\ell = E_{\rm target}$.

The results are mixed: when coherence decays steadily with depth, the \RCF\ schedule moves
the entangling angle onto the last layers and gains
$\Delta F = +0.022$ ($E_{\rm target} = 1$) and $+0.038$
($E_{\rm target} = 2$) over the rate-only schedule, most of what the
numerical optimum gains.
When the rotations make the affinity jump from layer to layer, the
schedule puts $\theta \approx 0.6$ on the single lowest-affinity
layer, well outside the small-angle regime in which
Eq.~\eqref{eq:linear-cost} holds, and loses $\Delta F = -0.035$.
Over 300 random circuits
($\gamma_\ell \in [0.02, 0.30]$, random rotations and idle
dephasing), the \RCF\ schedule beats the rate-only schedule on average
by $\Delta F = +0.003$ to $+0.005$ for
$E_{\rm target} \in \{0.5, 1, 2\}$ (standard errors 0.0005--0.0014),
but it is worse on 13--41\% of circuits and recovers only 10--25\% of
the gap to the numerical optimum, which is itself
$0.010$ to $0.052$.
Weighting each layer instead by the first-order infidelity
$\sum_{i\neq j} h_{ij}\,\lvert\rho_{ij}\rvert^2$ of the state exposed
to the gate noise, where $h_{ij}$ is the Hamming distance between
pointer states $i$ and $j$, recovers 23--38\%, although it fails in
the same way when the affinity jumps between layers.
That weight scales as $\bar\alpha^2$ rather than $\bar\alpha$, because
dephasing costs fidelity in proportion to the squared off-diagonal
elements.
The affinity is therefore operationally relevant, and the direction in
which it moves the schedule is right, but the linear-in-$\bar\alpha$
cost of Eq.~\eqref{eq:layer-cost} is not the best way to use it.

I tried the obvious repair: weight each layer by
$(\bar\alpha_\ell^{\rm pre})^2$ instead of $\bar\alpha_\ell^{\rm pre}$,
and add the bound $\theta_\ell \le \theta_{\max}$ to
Eq.~\eqref{eq:alloc-problem}.
The bound leaves Theorem~\ref{thm:alloc} intact; the optimal angle is
simply Eq.~\eqref{eq:alloc-solution} clipped at $\theta_{\max}$.
I chose $\theta_{\max} = \pi/12$ on a separate calibration set of 150
random circuits and then fixed it for the 300 reported above.
Of the two changes, the bound does most of the work.
In the case where the affinity jumps between layers, the capped
quadratic schedule gains $\Delta F = +0.019$ ($E_{\rm target} = 1$)
and $+0.027$ ($E_{\rm target} = 2$) where the linear schedule lost
$0.035$ and $0.039$; with steady decay it gives up a little
($+0.020$ and $+0.026$ instead of $+0.022$ and $+0.038$); and with
heterogeneous rates it gains $+0.011$ and $+0.012$.
Over the random circuits at $E_{\rm target} = 1$ it gains
$\Delta F = +0.013 \pm 0.002$, three times the linear schedule, recovers
43\% of the gap to the optimum instead of 14\%, and its worst single
loss falls from $0.11$ to $0.06$.
At $E_{\rm target} = 0.5$ the bound never binds and squaring the
weight adds little ($+0.003$).
At $E_{\rm target} = 2$ a fixed $\theta_{\max} = \pi/12$ is too
tight: meeting the target then forces angle onto at least four layers,
including noisy ones, and the gain falls to $+0.006$.
Squaring the weight without the bound barely changes the average and
makes the worst cases worse.
The lesson is that the coherence cost needs to stay inside the regime
in which it was derived, and that the bound should scale with
$E_{\rm target}/L$ rather than be fixed.
Even the best of these schedules recovers well under half of what
direct numerical optimization achieves, so a per-layer linear cost of
this kind, however it is weighted, remains an approximation.

\begin{figure*}[htbp]
\centering
\includegraphics[width=\textwidth]{fig4_alpha_varying}
\caption{
  Allocation with layer-varying pre-gate affinity ($L = 8$,
  $E_{\rm target} = 1$ in the first three panels).
  Gray bars: $\bar\alpha_\ell^{\rm pre}$ along the rate-only schedule.
  Lines: exchange angles $\theta_\ell$ for the rate-only schedule,
  the \RCF\ schedule ($\tilde\gamma_\ell \propto
  \bar\alpha_\ell^{\rm pre}\gamma_\ell$), the capped quadratic
  schedule ($\tilde\gamma_\ell \propto
  (\bar\alpha_\ell^{\rm pre})^2\gamma_\ell$,
  $\theta_\ell \le \pi/12$), and the numerical optimum, with final-state fidelities
  in the legends.
  \textbf{A:} uniform $\gamma = 0.15$, coherence decaying through idle
  dephasing.
  \textbf{B:} uniform $\gamma = 0.15$, interleaved single-qubit
  rotations.
  \textbf{C:} the heterogeneous rates of Fig.~\ref{fig:rcb-sweep}
  with idle dephasing.
  \textbf{Right:} mean fidelity gain over the rate-only schedule
  across 300 random circuits (error bars: standard error).
  Generated by \texttt{figures/\allowbreak benchmark\_\allowbreak alpha\_\allowbreak varying.py}.
}
\label{fig:alpha-varying}
\end{figure*}

\subsection{Scope of Theorem~\ref{thm:alloc}}
\label{subsec:alloc-scope}

The allocation theorem holds under the assumptions made explicit in
Sec.~\ref{sec:sim}: (i)~Markovian pure-dephasing dynamics in the
einselected pointer basis, with rates that are constant within each
gate window;
(ii)~gate-time-proportional coherence-limited error (small-$p$
leading-order linearization of Eq.~\eqref{eq:peff});
(iii)~direct exchange calibration so that $\URA(\theta)$ is realized
as a single tunable pulse rather than decomposed into fixed native
gates.
Where (iii) fails, the constant $\kappa$ relating $\theta$ to
duration in Eq.~\eqref{eq:layer-cost} is replaced by a
decomposition-depth--dependent factor; the optimization problem and
its solution form are unchanged but the numerical advantage shrinks
as in Sec.~\ref{sec:disc}.
Coherent calibration errors, leakage to non-computational levels, and
crosstalk are not modelled by Eq.~\eqref{eq:peff} and can shift the
optimum; their incorporation as additional per-layer cost terms is
straightforward and preserves the convex-program structure.
